\section{Manifolds of Constant Curvature}

Let $(M,g)$ and $(\widetilde M,\widetilde g)$ be Riemannian manifolds of the
same dimension, let $U\subseteq M$ be connected and open, and let
$\varphi:U\to\widetilde M$ be a local isometry.  For a continuous path
$\gamma:[0,1]\to M$ with $\gamma(0)\in U$, an \emph{analytic continuation of
$\varphi$ along $\gamma$} is a family $(U_t,\varphi_t)_{t\in[0,1]}$ such that
$U_t$ is a connected neighborhood of $\gamma(t)$, each
$\varphi_t:U_t\to\widetilde M$ is a local isometry, $\varphi_0=\varphi$ on
$U_0\cap U$, and, for every $t$, there is $\delta>0$ such that
$|t-t_1|<\delta$ implies $\gamma(t_1)\in U_t$ and $\varphi_t=\varphi_{t_1}$
on $U_t\cap U_{t_1}$.  No compatibility is required between distinct return
times that are not sufficiently close.

\begin{lemma}[Uniqueness of Analytic Continuation]
\label{lem:uniqueness-analytic-continuation}
\dcref{ch12:12.1}
With $(M,g)$, $(\widetilde M,\widetilde g)$, and $\varphi$ as above, two
analytic continuations of $\varphi$ along the same path $\gamma$ agree on a
neighborhood of $\gamma(1)$ at time $1$.
\begin{proof}
\sketch Let $\mathcal T$ be the set of times at which the two continuations
agree near $\gamma(t)$.  It contains $0$ and is open by the local compatibility
condition.  If $t_i\in\mathcal T$ and $t_i\to t$, compatibility identifies both
continuations near $\gamma(t_i)$ for all sufficiently large $i$.  Their values
and differentials therefore agree at $\gamma(t)$ by continuity, and the local
uniqueness theorem for local isometries implies agreement near $\gamma(t)$.
Thus $\mathcal T$ is also closed, so $\mathcal T=[0,1]$.
\end{proof}
\end{lemma}

\begin{theorem}[Monodromy Theorem for Local Isometries]
\label{thm:monodromy-theorem-local-isometries}
\dcref{ch12:12.2}
Let $(M,g)$ and $(\widetilde M,\widetilde g)$ be connected Riemannian
manifolds.  Suppose $U$ is a neighborhood of $p\in M$ and
$\varphi:U\to\widetilde M$ is a local isometry admitting analytic continuation
along every path starting at $p$.  If $\gamma_0$ and $\gamma_1$ are
path-homotopic paths from $p$ to $q$, their analytic continuations agree on a
neighborhood of $q$.
\begin{proof}
\sketch Let $H_s$ be a path homotopy from $\gamma_0$ to $\gamma_1$, with
fixed endpoints.  A finite subfamily of any continuation along $H_s$ covers
the compact path, and the same local isometries form a continuation along
$H_{s'}$ for all sufficiently close $s'$.  By uniqueness of continuation,
the endpoint value is therefore locally constant in $s$, hence constant.
Modifying the paths only in a fixed neighborhood of their common endpoint
gives the same conclusion at every nearby endpoint, proving local agreement.
\end{proof}
\end{theorem}

\begin{corollary}
\label{cor:global-isometry-from-local-simply-connected}
\uses{thm:monodromy-theorem-local-isometries, thm:completeness-of-covering-map, cor:coverings-simply-connected}
\dcref{ch12:12.3}

Let $(M,g)$ and $(\widetilde M,\widetilde g)$ be simply connected, complete
Riemannian manifolds.  If a local isometry $\varphi:U\to\widetilde M$ on a
connected neighborhood $U$ of $p\in M$ admits analytic continuation along
every path starting at $p$, then it extends to a global isometry
$\Phi:M\to\widetilde M$.
\begin{proof}
\sketch Define $\Phi(q)$ by continuing $\varphi$ along a path from $p$ to
$q$.  Simple connectedness of $M$ and the monodromy theorem make this a
well-defined local isometry extending $\varphi$.  Completeness of $M$ makes
$\Phi$ a Riemannian covering, and simple connectedness of $\widetilde M$ makes
that covering bijective.
\end{proof}
\end{corollary}

\begin{theorem}[Killing-Hopf]
\label{thm:killing-hopf}
\uses{cor:local-uniqueness-constant-curvature, cor:global-isometry-from-local-simply-connected, thm:small-geodesic-ball-convex, prop:model-manifolds-isometric-actions, prob:model-manifolds-isometric-actions}
\dcref{ch12:12.4}

Let $(M,g)$ be a complete, simply connected Riemannian $n$-manifold of
constant sectional curvature, where $n\geq2$.  Then $M$ is isometric to one
of $\mathbb R^n$, $S^n(R)$, or $\mathbb H^n(R)$.
\begin{proof}
\sketch Choose the simply connected model $(\widetilde M,\widetilde g)$ with
the same curvature.  Local uniqueness for constant-curvature metrics gives a
local isometry $\varphi$ from a neighborhood of any $p\in M$ into
$\widetilde M$.  Along a path from $p$, choose finitely many successive convex
geodesic balls $B_i$, each isometric to an open subset of the model.  Starting
with $\psi_0=\varphi|_{B_0}$, compose a model chart on $B_i$ with a global
model isometry extending $\psi_{i-1}$ on the connected overlap
$B_{i-1}\cap B_i$.  This inductively produces compatible local isometries and
hence an analytic continuation along the path.  The preceding corollary then
extends $\varphi$ to a global isometry.
\end{proof}
\end{theorem}

\begin{corollary}[Characterization of Constant-Curvature Manifolds]
\label{cor:characterization-constant-curvature}
\uses{thm:killing-hopf, prop:covering-automorphisms-discrete-lie-group, cor:isometric-to-quotient, prop:model-manifolds-isometric-actions, prop:discrete-subgroups-dimension-zero, prop:isometry-group-acts-properly, prop:riemannian-quotient-metric, cor:complete-iff-covering-complete, prob:model-manifolds-isometric-actions, prob:isometry-group-acts-properly}
\dcref{ch12:12.5}

Up to isometry, the complete, connected, constant-curvature Riemannian
$n$-manifolds are exactly the quotients $\widetilde M/\Gamma$, where
$\widetilde M$ is $\mathbb R^n$, $S^n(R)$, or $\mathbb H^n(R)$ and $\Gamma$
is a discrete subgroup of $\operatorname{Iso}(\widetilde M)$ acting freely.
\begin{proof}
\sketch For a complete constant-curvature manifold $M$, equip its universal
cover $\pi:\widetilde M\to M$ with the pullback metric.  It is complete and,
by Killing--Hopf, isometric to the appropriate model.  Its deck group
$\Gamma$ acts freely, properly, and isometrically, and $M$ is isometric to
$\widetilde M/\Gamma$.  The deck orbit of each point is closed and discrete;
continuity of the isometry action therefore rules out an accumulation point
of $\Gamma$ in $\operatorname{Iso}(\widetilde M)$.

Conversely, a discrete subgroup acting freely is a closed Lie subgroup and
acts properly on the model.  The quotient carries the unique metric making
the projection a Riemannian covering.  It is complete and has the model's
constant sectional curvature.
\end{proof}
\end{corollary}

A complete, connected constant-curvature Riemannian manifold is a
\emph{space form}; it is \emph{spherical}, \emph{Euclidean}, or
\emph{hyperbolic} according as its sectional curvature is positive, zero, or
negative.  Its classification reduces to classifying the fixed-point-free
discrete subgroups of $O(n+1)$, $E(n)$, or $O^+(n,1)$, respectively; see
\cite{Wol11}.

\section{Symmetric Spaces Revisited}

\begin{theorem}
\label{thm:parallel-riemann-curvature-implies-symmetry}
\uses{thm:characterization-locally-symmetric-spaces, cor:global-isometry-from-local-simply-connected}
\dcref{ch12:12.6}

Every complete, simply connected Riemannian manifold with parallel Riemann
curvature tensor is a symmetric space.
\begin{proof}
\sketch Fix $p\in M$.  The local characterization of symmetric spaces gives
a point reflection $\varphi$ near $p$.  Along any path from $p$, cover the
path by finitely many successive convex geodesic balls.  If a local isometry
has been constructed on one ball, its value and differential at a point of
the next overlap give a linear isometry preserving the metric and curvature.
The local construction theorem for parallel curvature tensors produces the
next local isometry, and local uniqueness makes the two maps agree on the
connected overlap.  Thus $\varphi$ admits analytic continuation along every
path and extends globally.  The extension has differential
$-\operatorname{Id}$ at $p$, so it is a global point reflection.
\end{proof}
\end{theorem}

\begin{corollary}
\label{cor:symmetric-space-equivalence}
\uses{thm:characterization-locally-symmetric-spaces, thm:parallel-riemann-curvature-implies-symmetry, cor:complete-iff-covering-complete, prop:covering-automorphisms-discrete-lie-group, cor:isometric-to-quotient, prop:riemannian-quotient-metric, prop:symmetric-space-parallel-curvature, prob:symmetric-space-parallel-curvature}
\dcref{ch12:12.7}

For a complete, connected Riemannian manifold $(M,g)$, the following are
equivalent:
\begin{enumerate}
\item[(a)] The curvature tensor of $M$ is parallel.
\item[(b)] The manifold $M$ is locally symmetric.
\item[(c)] The manifold $M$ is isometric to $\widetilde M/\Gamma$, where
$\widetilde M$ is a globally symmetric space and $\Gamma$ is a discrete Lie
group acting freely, properly, and isometrically on $\widetilde M$.
\end{enumerate}
\begin{proof}
\sketch The local characterization gives (a)$\Leftrightarrow$(b).  Under (a),
the universal cover with its pullback metric is complete, simply connected,
and has parallel curvature, hence is globally symmetric by the preceding
theorem.  Its deck group has the properties in (c), and the deck quotient is
$M$.  Conversely, a globally symmetric space has parallel curvature, which
descends through the Riemannian covering $\widetilde M\to\widetilde M/\Gamma$.
\end{proof}
\end{corollary}

\section{Manifolds of Nonpositive Curvature}

\begin{theorem}[Cartan--Hadamard]
\label{thm:cartan-hadamard}
\lean{LeeLib.Ch12.operator_nonpos_of_sectional_nonpos, LeeLib.Ch12.cartanHadamardModel_simplyConnected, LeeLib.Ch12.cartanHadamardModelDiffeomorph, LeeLib.Ch12.cartanHadamardExpMap_sectional, LeeLib.Ch12.cartanHadamardCovering_sectional, LeeLib.Ch12.cartanHadamardDiffeomorph_sectional, LeeLib.Ch12.cartanHadamardDiffeomorph_euclidean_sectional, LeeLib.Ch12.cartanHadamardDiffeomorph_sectional_coe}
\uses{thm:conjugate-point-comparison-I, prop:critical-points-conjugate, cor:exponential-map-defined-implies-complete, thm:completeness-of-covering-map}
\leanok
\dcref{ch12:12.8}



If $(M,g)$ is a complete, connected Riemannian manifold with nonpositive
sectional curvature, then $\exp_p:T_pM\to M$ is a smooth covering map for
every $p\in M$.  Consequently, the universal cover of $M$ is diffeomorphic to
$\mathbb R^n$, and a simply connected $M$ is itself diffeomorphic to
$\mathbb R^n$.
\begin{proof}
\sketch Nonpositive curvature gives no conjugate points, so $\exp_p$ is a
local diffeomorphism on all of $T_pM$.  The pullback $\bar g=\exp_p^*g$ is
therefore a Riemannian metric for which $\exp_p$ is a local isometry.  The
radial lines through the origin are defined for all time as $\bar g$-geodesics,
so $(T_pM,\bar g)$ is complete.  A local isometry from a complete manifold is
a smooth covering.  Since $T_pM\cong\mathbb R^n$ is simply connected, it is
the universal cover, and the remaining claims follow.
\end{proof}
\end{theorem}

A \emph{Cartan--Hadamard manifold} is a complete, simply connected Riemannian
manifold of nonpositive sectional curvature.  A \emph{line} is the image of a
nonconstant geodesic defined on all of $\mathbb R$ whose restriction between
any two of its points is minimizing.

\begin{proposition}[Basic Properties of Cartan--Hadamard Manifolds]
\label{prop:basic-properties-cartan-hadamard}
\uses{lem:radial-distance-well-defined, prop:radial-geodesic-unique-minimizer, cor:radial-distance-riemannian-distance, cor:geodesic-ball-metric-ball}
\dcref{ch12:12.9}

If $(M,g)$ is a Cartan--Hadamard manifold, then:
\begin{enumerate}[label=(\alph*)]
\item the injectivity radius of $M$ is infinite;
\item the image of every nonconstant maximal geodesic is a line;
\item every two distinct points lie on a unique line;
\item every open or closed metric ball is a geodesic ball;
\item for each $q\in M$, the function $r(x)=d_g(q,x)$ is smooth on
$M\setminus\{q\}$ and $r^2$ is smooth on $M$.
\end{enumerate}
\begin{proof}
\sketch Cartan--Hadamard makes $M$ a normal neighborhood of each of its
points.  The stated properties follow from the uniqueness and minimizing
properties of radial geodesics, the identification of radial and Riemannian
distance, and the corresponding description of geodesic balls.
\end{proof}
\end{proposition}

Three points in a Cartan--Hadamard manifold are \emph{noncollinear} if they do
not lie on one line.  Their \emph{geodesic triangle} $\triangle ABC$ is the
union of the three unique connecting geodesic segments.  The angle $\angle A$
is the angle between the initial velocities from $A$ to $B$ and from $A$ to
$C$, and similarly at the other vertices.

\begin{proposition}
\label{prop:geodesic-triangle-inequalities}
\dcref{ch12:12.10}
Let $\triangle ABC$ be a geodesic triangle in a Cartan--Hadamard manifold,
and let $a,b,c$ be the side lengths opposite $A,B,C$, respectively.  Then
\begin{enumerate}
\item[(a)] $c^2\geq a^2+b^2-2ab\cos\angle C$;
\item[(b)] $\angle A+\angle B+\angle C\leq\pi$.
\end{enumerate}
If all sectional curvatures are strictly negative, both inequalities are
strict.
\end{proposition}

\begin{corollary}
\label{cor:no-simply-connected-compact-nonpositive}
\dcref{ch12:12.11}
No simply connected compact manifold admits a Riemannian metric of
nonpositive sectional curvature.
\end{corollary}

\begin{corollary}
\label{cor:product-manifolds-nonpositive}
\uses{cor:simply-connected-orientable}
\dcref{ch12:12.12}

Let $M$ and $N$ be positive-dimensional compact, connected smooth manifolds,
at least one of which is simply connected.  Then $M\times N$ admits no
Riemannian metric of nonpositive sectional curvature.
\begin{proof}
\sketch Suppose, after interchanging the factors, that $M$ is simply connected
and that $M\times N$ has such a metric.  For the universal cover
$\widetilde N\to N$, the induced metric on $M\times\widetilde N$ is complete
and nonpositively curved, so Cartan--Hadamard makes this product diffeomorphic
to Euclidean space.  The compact simply connected manifold $M$ is orientable;
let $\mu$ be a positive orientation form.  It is closed and nonexact because
$\int_M\mu>0$.  If $p:M\times\widetilde N\to M$ is projection, then $p^*\mu$
is exact because the product is contractible.  For
$\sigma(x)=(x,y_0)$, however,
\[
\int_{\sigma(M)}p^*\mu=\int_M\sigma^*p^*\mu=\int_M\mu>0,
\]
contradicting the vanishing of the integral of an exact form over the closed
submanifold $\sigma(M)$.
\end{proof}
\end{corollary}

\begin{example}
\label{ex:positive-curvature-on-product-spheres}
\dcref{ch12:12.13}
For $m\geq2$ and $n\geq1$, every Riemannian metric on $S^m\times S^n$ has
positive sectional curvature at some point.
\end{example}

A space is \emph{aspherical} if $\pi_k$ vanishes for every $k>1$.

\begin{corollary}
\label{cor:nonpositive-implies-aspherical}
\dcref{ch12:12.14}
If a smooth manifold $M$ admits a complete Riemannian metric of nonpositive
sectional curvature, then $M$ is aspherical.
\begin{proof}
\sketch Cartan--Hadamard identifies the universal cover with $\mathbb R^n$,
which is contractible.  Covering maps induce isomorphisms on $\pi_k$ for
$k>1$, so all higher homotopy groups of $M$ vanish.
\end{proof}
\end{corollary}

\section*{Cartan's Torsion Theorem}

\begin{lemma}
\label{lem:geodesic-convexity-distance-squared}
\uses{prop:basic-properties-cartan-hadamard}
\dcref{ch12:12.15}

Let $(M,g)$ be a Cartan--Hadamard manifold, fix $q\in M$, and set
$f(x)=\frac12d_g(x,q)^2$.  For every nonconstant geodesic segment
$\gamma:[0,1]\to M$ and every $t\in(0,1)$,
\begin{equation}
f(\gamma(t))<(1-t)f(\gamma(0))+tf(\gamma(1)). \tag{12.1}
\end{equation}
Thus $f$ is strictly geodesically convex.
\begin{proof}
\sketch Write $f=\frac12r^2$, where $r$ is radial distance from $q$.
The basic properties above make $f$ smooth on $M$, and Hessian comparison
gives, away from $q$,
\[
\mathcal H_f=\operatorname{grad}r\otimes dr+r\mathcal H_r
\geq\operatorname{grad}r\otimes dr+\pi_r=\operatorname{Id}.
\]
Continuity extends the inequality to $q$.  Along a nonconstant geodesic,
$(f\circ\gamma)''=\langle\mathcal H_f(\gamma'),\gamma'\rangle
\geq|\gamma'|^2>0$, which yields the displayed strict convexity inequality.
\end{proof}
\end{lemma}

\begin{lemma}
\label{lem:unique-minimal-enclosing-ball}
\uses{lem:geodesic-convexity-distance-squared}
\dcref{ch12:12.16}

Let $(M,g)$ be a Cartan--Hadamard manifold and let $S\subseteq M$ be compact
with more than one point.  Then there is a unique closed ball of minimum
radius containing $S$.
\begin{proof}
\sketch Put $\delta=\operatorname{diam}S>0$, choose $q_0\in S$, and let $c_0$
be the infimum of radii of closed balls containing $S$.  For a minimizing
sequence $\overline B_{c_i}(p_i)$ with $c_i\searrow c_0$, eventually
$p_i\in\overline B_{2\delta}(q_0)$.  This ball is compact, so a subsequence of
the centers converges to $p_0$.  The triangle inequality then gives
$S\subseteq\overline B_{c_0}(p_0)$.

If a second minimum ball has a distinct center $p_0'$, let $m$ be the midpoint
of the geodesic segment joining the centers.  For each $q\in S$, strict
convexity of $x\mapsto\frac12d_g(q,x)^2$ gives
\[
\frac12d_g(q,m)^2<\frac14d_g(q,p_0)^2+
\frac14d_g(q,p_0')^2\leq\frac12c_0^2.
\]
Compactness of $S$ makes $\max_{q\in S}d_g(q,m)<c_0$, contradicting the
minimality of $c_0$.
\end{proof}
\end{lemma}

The center of this minimum enclosing ball is the \emph{1-center} of $S$.

\begin{theorem}[Cartan's Fixed-Point Theorem]
\label{thm:cartans-fixed-point}
\uses{lem:unique-minimal-enclosing-ball}
\dcref{ch12:12.17}

Let $(M,g)$ be a Cartan--Hadamard manifold and let a compact Lie group $G$
act smoothly and isometrically on $M$.  Then some $p_0\in M$ is fixed by
every element of $G$.
\begin{proof}
\sketch For $q_0\in M$, the orbit $S=G\cdot q_0$ is compact.  If it is a
singleton, its point is fixed.  Otherwise, let $p_0$ be its unique 1-center.
Every $\varphi\in G$ preserves $S$ and sends its minimum enclosing ball to a
ball of the same radius centered at $\varphi\cdot p_0$.  Uniqueness of the
1-center implies $\varphi\cdot p_0=p_0$.
\end{proof}
\end{theorem}

An element $\varphi$ of a group is a \emph{torsion element} if
$\varphi^k=1$ for some integer $k>1$.  A group is \emph{torsion-free} if its
only torsion element is the identity.

\begin{corollary}[Cartan's Torsion Theorem]
\label{cor:cartans-torsion-theorem}
\uses{thm:cartans-fixed-point}
\dcref{ch12:12.18}

If $(M,g)$ is a complete, connected Riemannian manifold of nonpositive
sectional curvature, then $\pi_1(M)$ is torsion-free.
\begin{proof}
\sketch The universal cover $(\widetilde M,\widetilde g)$ is a
Cartan--Hadamard manifold, and its deck group
$\Gamma\cong\pi_1(M)$ acts freely and isometrically.  A torsion element
$\varphi\in\Gamma$ generates a finite, hence compact, Lie group.  Cartan's
fixed-point theorem gives this group a fixed point, while freeness forces
$\varphi$ to be the identity.
\end{proof}
\end{corollary}
\section{Preissman's Theorem}

\begin{theorem}[Preissman]
\label{thm:preissman-theorem}
\dcref{ch12:12.19}
If $(M,g)$ is compact, connected, and has strictly negative sectional
curvature, then every nontrivial abelian subgroup of $\pi_1(M)$ is isomorphic
to $\mathbb Z$.
\end{theorem}

For a complete Riemannian manifold, an isometry $\varphi$ has an \emph{axis} if
some geodesic $\gamma:\mathbb R\to M$ satisfies
$\varphi(\gamma(t))=\gamma(t+c)$ for a nonzero constant $c$.  An isometry with
an axis and no fixed point is \emph{axial}.

\begin{example}[Axial Isometries]
\label{ex:axial-isometries}
\dcref{ch12:12.20}
Every nontrivial Euclidean translation is axial, with axes parallel to its
translation vector.  In the upper half-plane model, a dilation
$(x,y)\mapsto(cx,cy)$ with $c\ne1$ is axial with axis
$\gamma(t)=(0,e^t)$; nontrivial horizontal translations are not axial.
\end{example}

\begin{lemma}[Axes of Covering Automorphisms]
\label{lem:covering-automorphism-axis}
\dcref{ch12:12.21}
If $(M,g)$ is compact and connected, and
$\pi:(\widetilde M,\widetilde g)\to(M,g)$ is its universal Riemannian cover,
then every nonidentity covering automorphism has an axis.  The projected axis
is a closed geodesic shortest in its free homotopy class.
\end{lemma}

\begin{proof}
\sketch
Fix a nonidentity deck transformation $\varphi$.  Paths in $\widetilde M$
from $\widetilde x$ to $\varphi(\widetilde x)$ project to one free homotopy
class of loops in $M$; path lifting and monodromy show conversely that every
loop in that class has such a lift.  The class is nontrivial because $\varphi$
has no fixed point.  Compactness gives a shortest closed geodesic $\gamma$ in
the class.  A lift $\widetilde\gamma$ satisfies
$\widetilde\gamma(1)=\varphi(\widetilde\gamma(0))$ and extends globally by
completeness.  The lifted geodesics
$t\mapsto\varphi(\widetilde\gamma(t))$ and
$t\mapsto\widetilde\gamma(t+1)$ have the same initial point and projected
initial velocity; injectivity of $d\pi$ and geodesic uniqueness make them
equal, so $\widetilde\gamma$ is an axis.
\end{proof}

\begin{lemma}[Uniqueness of Axes in Negative Curvature]
\label{lem:unique-axis-cartan-hadamard}
\uses{prop:geodesic-triangle-inequalities, prop:round-metric-distance-sphere, prob:round-metric-distance-sphere}
\dcref{ch12:12.22}

If $(M,g)$ is Cartan--Hadamard with strictly negative sectional curvature, every
axial isometry has a unique axis up to reparametrization.
\end{lemma}

\begin{proof}
\sketch
Normalize two putative axes so that the isometry translates each by time $1$.
If they are disjoint, join their initial points by a geodesic segment and its
translate; these form a geodesic quadrilateral.  Isometry invariance makes the
quadrilateral angle sum $2\pi$, while its two negative-curvature triangular
decompositions have angle sums strictly less than $\pi$.  The angle triangle
inequality (equivalently, the round metric triangle inequality on a unit
tangent sphere) makes the quadrilateral angle sum strictly less than $2\pi$,
a contradiction.  If the axes intersect, their translates give a second
intersection and Cartan--Hadamard geodesic uniqueness forces coincidence up to
reparametrization.
\end{proof}

\begin{proof}
\sketch
Let $\pi:\widetilde M\to M$ be the universal cover of the compact negatively
curved manifold.  It is Cartan--Hadamard.  For a nontrivial abelian subgroup
$H$ of deck transformations, choose $\varphi\in H\setminus\{1\}$ and an axis
$\gamma$ for it.  Every $\psi\in H$ commutes with $\varphi$, so $\psi\circ\gamma$
is another axis; uniqueness gives a translation
\[
\psi(\gamma(t))=\gamma(t+F(\psi)).
\]
Backward reparametrization would give a fixed point for a nonidentity deck
transformation, so all translations preserve orientation.  The map
$F:H\to\mathbb R$ is an injective homomorphism.  If positive elements of
$F(H)$ had infimum $0$, or if the positive infimum were not attained, a
nontrivial element would translate the axis by less than
$\operatorname{inj}(M)$, contradicting injectivity radius.  The least positive
element therefore generates $F(H)$, so $H\cong\mathbb Z$.
\end{proof}

\begin{corollary}
\label{cor:no-product-negative-sectional-curvature}
\dcref{ch12:12.23}
No product of positive-dimensional connected compact manifolds admits a
metric with strictly negative sectional curvature.
\end{corollary}

\section{Manifolds of Positive Curvature}

\begin{theorem}[Myers Diameter and Compactness]
\label{thm:myers-positive-ricci-geometric}
\lean{LeeLib.Ch12.myersPositiveRicci_diameterCompact_of_dim}


Let $(M,g)$ be complete and connected, of dimension $n\geq2$.  If $R>0$ and
\[
\operatorname{Rc}(v,v)\geq\frac{n-1}{R^2}
\]
for every unit vector $v$, then $M$ is compact and
$\operatorname{diam}(M)\leq\pi R$.
\end{theorem}

\begin{proof}
\leanok
\sketch
The sine-field index inequality excludes minimizing geodesic segments longer
than $\pi R$.  Hopf--Rinow supplies minimizing segments between points, giving
the diameter bound; completeness and boundedness imply compactness.
\end{proof}

\begin{theorem}[Myers on an Explicit Simply Connected Riemannian Cover]
\label{thm:myers-explicit-simply-connected-cover}
\lean{LeeLib.Ch12.myers_of_simplyConnected_riemannianCover_of_complete_base}
\uses{thm:myers-positive-ricci-geometric, cor:complete-iff-covering-complete}
\leanok



Let $\pi:(\widetilde M,\widetilde g)\to(M,g)$ be a Riemannian cover of
connected $n$-manifolds, with $\widetilde M$ simply connected and $M$ complete.
If $n\geq2$, $R>0$, and
$\operatorname{Rc}_{\widetilde g}(v,v)\geq(n-1)/R^2$ for unit $v$, then $M$ is
compact, $\operatorname{diam}(M)\leq\pi R$, and $\pi_1(M,p)$ is finite.
\end{theorem}

\begin{proof}
\leanok
\sketch
By \cref{cor:complete-iff-covering-complete}, $\widetilde M$ is complete;
applying \cref{thm:myers-positive-ricci-geometric} makes it compact with
diameter at most $\pi R$.  The covering map is a
surjective local isometry and transfers the diameter bound and compactness to
$M$.  Monodromy identifies each fundamental group with a covering fiber; the
fiber is closed and discrete in compact $\widetilde M$, hence finite.
\end{proof}

\begin{theorem}[Myers]
\label{thm:myers-positive-ricci}
\uses{thm:myers-positive-ricci-geometric, thm:myers-explicit-simply-connected-cover, cor:complete-iff-covering-complete, prop:fiber-cardinality-fundamental-group}
\notready
\dcref{ch12:12.24}
\notready

Let $(M,g)$ be complete and connected, and suppose
$\operatorname{Rc}(v,v)\geq(n-1)/R^2$ for every unit vector, with $R>0$.  Then
$M$ is compact, has diameter at most $\pi R$, and has finite fundamental group.
\end{theorem}

\begin{proof}
\sketch
The diameter comparison theorem gives $\operatorname{diam}(M)\leq\pi R$.
The exponential map at a base point is therefore surjective on the compact
closed ball $\overline B_{\pi R}(0)$, proving compactness.  The universal cover
inherits completeness and the Ricci bound, so it is compact by the same
argument.  Its fibers are closed discrete subsets and correspond to
$\pi_1(M,p)$, hence are finite.
\end{proof}

\begin{corollary}
\label{cor:positive-ricci-finite-fundamental-group}
\uses{prop:unit-tangent-bundle-properties}
\dcref{ch12:12.25}

If $(M,g)$ is compact and connected and its Ricci tensor is positive definite,
then $\pi_1(M)$ is finite.
\end{corollary}

\begin{proof}
\sketch
Compactness of the unit tangent bundle gives a uniform $c>0$ with
$\operatorname{Rc}(v,v)\geq c$ for unit $v$.  Apply Myers with
$c=(n-1)/R^2$.
\end{proof}

\begin{corollary}
\label{cor:einstein-positive-scalar-curvature-compact}
\dcref{ch12:12.27}
Every complete Einstein manifold with positive scalar curvature is compact.
\end{corollary}

\begin{proof}
\sketch
For an Einstein metric, $\operatorname{Rc}=(S/n)g$.  Positive constant $S$
therefore satisfies Myers with $(n-1)/R^2=S/n$.
\end{proof}

The maximal-diameter rigidity theorem is due to Cheng~\cite{Che75}.

\begin{theorem}[Cheng's Maximal Diameter Theorem]
\label{thm:chengs-maximal-diameter}
\uses{cor:einstein-positive-scalar-curvature-compact}
\dcref{ch12:12.28}

If $(M,g)$ satisfies Myers's hypotheses and
$\operatorname{diam}(M)=\pi R$, then $(M,g)$ is isometric to
$(S^n(R),\bar g_R)$.
\end{theorem}

\begin{proof}
\sketch
Compactness gives $p_1,p_2$ with $d(p_1,p_2)=\pi R$.  Bishop--Gromov applied
to balls around each point gives, for $0<\delta<\pi R$,
\begin{align}
\frac{\operatorname{Vol}(B_\delta(p_i))}{\operatorname{Vol}(M)}
&\geq\frac{\operatorname{Vol}(\widetilde B_\delta)}
{\operatorname{Vol}(S^n(R))}. \label{eq:12.2}
\end{align}
If $\delta_1+\delta_2=\pi R$, the two balls are disjoint, so
\begin{align}
\operatorname{Vol}(B_{\delta_1}(p_1))
+\operatorname{Vol}(B_{\delta_2}(p_2))&\leq\operatorname{Vol}(M). \label{eq:12.3}
\end{align}
The corresponding spherical balls partition the sphere, forcing equality in
both inequalities and, since ball boundaries have measure zero,
\begin{align}
\operatorname{Vol}(\overline B_{\delta_1}(p_1))
+\operatorname{Vol}(\overline B_{\delta_2}(p_2))&=\operatorname{Vol}(M). \label{eq:12.4}
\end{align}
If $d(p_1,q)+d(p_2,q)>\pi R$, a neighborhood of $q$ avoids both closed balls,
contradicting (12.4).  Hence the distance sum is constantly $\pi R$.
Every geodesic from $p_1$ is minimizing to time $\pi R$, so the injectivity
radius at $p_1$ is $\pi R$ and $M\setminus\{p_2\}$ is a geodesic ball; similarly
for $p_2$.  The two distance functions are smooth off the poles and satisfy
$r_1+r_2=\pi R$.  Laplacian comparison with the spherical model (using
$c=1/R^2$) is therefore equality, forcing constant sectional curvature
$1/R^2$.  The constant-curvature covering is one-sheeted by the equal-volume
and injectivity-domain argument, so $M\cong S^n(R)$ isometrically.
\end{proof}

\begin{theorem}[Milnor]
\label{thm:milnor-polynomial-growth}
\dcref{ch12:12.29}
If $(M,g)$ is complete and connected with nonnegative Ricci curvature, every
finitely generated subgroup of $\pi_1(M)$ has polynomial growth of order at
most $n$.
\end{theorem}

For a finitely generated group $\Gamma$ with finite generating set $S$,
polynomial growth means that the number of elements representable by at most
$m$ generators and inverses is at most $Cm^k$ for some $C,k>0$ and all positive
integers $m$; its order is the infimum of such $k$.

\begin{corollary}
\label{cor:free-fund-group-no-ric-metric}
\uses{thm:milnor-polynomial-growth}
\dcref{ch12:12.30}

If a connected smooth manifold has fundamental group free on more than one
generator, it admits no complete metric with nonnegative Ricci curvature.
\end{corollary}

\begin{theorem}[Synge]
\label{thm:synge-positive-sectional}
\dcref{ch12:12.32}
Let $(M,g)$ be compact, connected, and strictly positively curved.
\begin{enumerate}
\item[(a)] If $n$ is even and $M$ is orientable, then $M$ is simply connected.
\item[(b)] If $n$ is odd, then $M$ is orientable.
\end{enumerate}
\end{theorem}

\begin{proof}
\sketch
Assume the relevant conclusion fails and pass to the universal cover
$\pi:\widetilde M\to M$.  Give $\widetilde M$ the pullback orientation in
the even case, and choose an orientation in the odd case.  A nontrivial deck
transformation $\varphi$ can be chosen orientation-preserving in the even case
and orientation-reversing in the odd case.  By the axis lemma, its axis projects
to a shortest closed geodesic; normalize the translation to time $1$.

Let $\bar P$ be parallel transport along the lifted geodesic and $P$ parallel
transport around its projection.  The covering diagram shows that $P$ is
orientation-preserving in the even case and orientation-reversing in the odd
case, hence $\det P=(-1)^n$.  Since $\gamma'(0)$ is a $+1$ eigenvector, the
eigenvalue parity forces a second orthogonal $+1$ eigenvector $v$.  Its parallel
extension $V$ is normal and periodic.

The variation $\Gamma(s,t)=\exp_{\gamma(t)}(sV(t))$ is a variation through
admissible loops.  The first variation vanishes, while the second variation is
\[
\left.\frac{d^2}{ds^2}\right|_{s=0}L(\Gamma_s)
=\int_0^1\bigl(|D_tV|^2-\operatorname{Rm}(V,\gamma',\gamma',V)\bigr)\,dt<0,
\]
because $V$ is parallel and sectional curvature is positive.  This contradicts
minimality of the closed geodesic in its free homotopy class.
\end{proof}

\begin{corollary}
\label{cor:fundamental-groups-even-dimension}
\uses{thm:synge-positive-sectional}
\dcref{ch12:12.33}

Let $(M,g)$ be compact, connected, even-dimensional, and strictly positively
curved.  If $M$ is orientable then $\pi_1(M)=0$; otherwise
$\pi_1(M)\cong\mathbb Z/2$.
\end{corollary}

\begin{proof}
\sketch
The orientable case is Synge's theorem.  In the nonorientable case, the
two-sheeted orientable cover is simply connected by Synge, hence universal;
its fibers have two points, so the fundamental group is $\mathbb Z/2$.
\end{proof}
\section{Further Results}

For $\delta>0$, a Riemannian manifold is \emph{$\delta$-pinched} if some
constant $c>0$ satisfies
\[
\delta c\leq\sec(\Pi)\leq c
\]
for every tangent $2$-plane $\Pi$. It is \emph{strictly
$\delta$-pinched} if at least one inequality is strict.

\begin{theorem}[The Sphere Theorem]
\label{thm:sphere-theorem}
\dcref{ch12:12.34}

Every complete, simply connected Riemannian $n$-manifold that is strictly
$\frac14$-pinched is homeomorphic to $S^n$.
\end{theorem}

This theorem is due to Berger and Klingenberg [Ber60, Kli61]. Its pinching
constant is sharp in even dimensions because the Fubini--Study metrics are
$\frac14$-pinched by \cref{prop:complex-projective-space-curvature}.

\begin{theorem}[The Soul Theorem]
\label{thm:soul-theorem}
\uses{prop:complex-projective-space-curvature, prob:complex-projective-space-curvature}
\dcref{ch12:12.35}


Let $(M,g)$ be complete, connected, noncompact, and have nonnegative
sectional curvature. There is a compact totally geodesic submanifold
$S\subseteq M$, called a \emph{soul}, such that $M$ is diffeomorphic to the
normal bundle of $S$ in $M$. If the sectional curvature is strictly positive,
then $S$ is a point and $M$ is diffeomorphic to Euclidean space.
\end{theorem}

This is the Cheeger--Gromoll soul theorem [CG72].

A \emph{line} in a Riemannian manifold is the image of a nonconstant geodesic
$\gamma:\mathbb R\to M$ whose restriction to every compact interval is
minimizing.

\begin{theorem}[The Splitting Theorem]
\label{thm:splitting-theorem}
\dcref{ch12:12.36}

Let $(M,g)$ be complete, connected, noncompact, and have nonnegative Ricci
curvature. If $M$ contains a line, then there is a Riemannian manifold
$(N,h)$ with nonnegative Ricci curvature such that
\[
(M,g)\cong(N\times\mathbb R,h\oplus dt^2)
\]
isometrically.
\end{theorem}

This is the Cheeger--Gromoll splitting theorem [CG71]. Proofs of the sphere,
soul, and splitting theorems are given in [Pet16].

\begin{theorem}[Hamilton's 3-Manifold Theorem]
\label{thm:hamilton-3-manifold}
\dcref{ch12:12.37}

If $M$ is a simply connected compact Riemannian $3$-manifold with
positive-definite Ricci curvature, then $M$ is diffeomorphic to $S^3$.
\end{theorem}

This is Hamilton's theorem [Ham82].

\begin{theorem}[Hamilton's 4-Manifold Theorem]
\label{thm:hamilton-4-manifold}
\uses{prop:curvature-operator, prob:curvature-operator}
\dcref{ch12:12.38}


If $M$ is a simply connected compact Riemannian $4$-manifold with positive
curvature operator, then $M$ is diffeomorphic to $S^4$.
\end{theorem}

This is Hamilton's theorem [Ham86].

The all-dimensional extension is due to Böhm and Wilking \cite{BW08}.

\begin{theorem}[Böhm--Wilking]
\label{thm:bohm-wilking}
\uses{thm:killing-hopf}
\dcref{ch12:12.39}


Every simply connected compact Riemannian manifold with positive curvature
operator is diffeomorphic to a sphere.
\end{theorem}

The Hamilton and Böhm--Wilking results use the \emph{Ricci flow}
\[
\frac{\partial}{\partial t}g_t=-2\mathrm{Rc}_t.
\]
Under their respective curvature hypotheses, a rescaling of $g_t$ converges
to a metric of constant positive sectional curvature; the Killing--Hopf
theorem then gives the stated diffeomorphism classification.

A Riemannian metric is \emph{pointwise $\delta$-pinched} if for every
$p\in M$ there is a number $c(p)>0$ such that
\[
\delta c(p)\leq\sec(\Pi)\leq c(p)
\qquad(\Pi\subseteq T_pM).
\]
It is \emph{strictly pointwise $\delta$-pinched} if at each point at least
one inequality is strict. Brendle and Schoen proved the following refinement
\cite{BS09}.

\begin{theorem}[Differentiable Sphere Theorem]
\label{thm:differentiable-sphere-theorem}
\dcref{ch12:12.40}

Let $(M,g)$ be compact and simply connected with dimension $n\geq4$. If $g$
is strictly pointwise $\frac14$-pinched, then $M$ is diffeomorphic to $S^n$
with its standard smooth structure.
\end{theorem}

See \cite{Bre10} for a proof.

\begin{exercise}
\label{exer:paraboloid-positive-curvature}
\dcref{ch12:12.26}
\uses{prob:graph-submanifold-shape}
Let $n \geq 2$, and let $M \subseteq \mathbb{R}^{n+1}$ be the paraboloid $\{(x^1, \ldots, x^n, y) : y = |x|^2\}$ with the induced metric (see Problem 8-1). Show that $M$ has strictly positive sectional and Ricci curvatures everywhere, but is not compact.
\end{exercise}

\begin{exercise}
\label{exer:prove-corollary-12-30}
\dcref{ch12:12.31}
\uses{cor:free-fund-group-no-ric-metric}
Prove this corollary.
\end{exercise}

\begin{exercise}
\label{prob:cartan-hadamard-generalization}
\dcref{ch12:12-9}
Prove the following generalization of the Cartan--Hadamard theorem, due to Robert Hermann [Her63]: Suppose $(M,g)$ is a complete, connected Riemannian manifold with nonpositive sectional curvature, and $P \subseteq M$ is a connected, properly embedded, totally geodesic submanifold. If for some $x \in P$, the homomorphism $\pi_1(P,x) \to \pi_1(M,x)$ induced by the inclusion $P \hookrightarrow M$ is surjective, then the normal exponential map $E : NP \to M$ is a diffeomorphism.
\end{exercise}

\begin{exercise}[Classification of Flat Tori]
\label{prob:classification-flat-tori}
\dcref{ch12:12-5}
Let $T^2 = \mathbb{S}^1 \times \mathbb{S}^1$ denote the 2-torus. Show that if $g$ is a flat Riemannian metric on $T^2$, then $(T^2, g)$ is isometric to one and only one Riemannian quotient $\mathbb{R}^2/\Lambda$, where $\Lambda$ is a lattice generated by a basis of the form $((a,0), (b,c))$, with $a > 0$, $0 \leq b \leq a/2$, $c > 0$, and $b^2 + c^2 \geq a^2$. [Hint: Given a lattice $\Lambda \subseteq \mathbb{R}^2$, let $v_1$ be an element of $\Lambda \setminus \{(0,0)\}$ of minimal norm; let $v_2$ be an element of $\Lambda \setminus \langle v_1 \rangle$ of minimal norm (where $\langle v_1 \rangle$ is the cyclic subgroup generated by $v_1$), chosen so that the angle between $v_1$ and $v_2$ is less than or equal to $\pi/2$; and then apply a suitable orthogonal transformation.]
\end{exercise}

\begin{exercise}
\label{prob:compact-euclidean-space-forms}
\dcref{ch12:12-3}
\uses{ex:open-mobius-band}
Show that every compact 2-dimensional Euclidean space form is diffeomorphic to the torus or the Klein bottle, and every noncompact one is diffeomorphic to $\mathbb{R}^2$, $\mathbb{S}^1 \times \mathbb{R}^2$, or the open Möbius band (see Example 2.35).
\end{exercise}

\begin{exercise}\label{prob:complete-lines-sectional-curvature}\dcref{ch12:12-13}
Suppose $(M, g)$ is a complete Riemannian manifold. Recall that a \emph{line} in $M$ is the image of a nonconstant geodesic defined on all of $\mathbb{R}$ that is minimizing between every pair of its points.

\begin{enumerate}[label=(\alph*)]
\item Show that if $(M, g)$ has strictly positive sectional curvature, then $M$ contains no lines. [Hint: Given a geodesic $\gamma : \mathbb{R} \to M$, let $\alpha : \mathbb{R} \to \mathbb{R}$ be a smooth positive function such that $\alpha(t)$ is smaller than the minimum sectional curvature at $\gamma(t)$ for each $t$, and let $u : \mathbb{R} \to \mathbb{R}$ be the solution to the initial value problem $u'' + \alpha u = 0$ with $u(0) = 1$ and $u'(0) = 0$. Prove that there are numbers $a < 0 < b$ such that $u(a) = u(b) = 0$, and evaluate $I(V, V)$ for a vector field of the form $V(t) = u(t)E(t)$ along $\gamma|_{[a,b]}$, where $E$ is a parallel unit normal vector field.]
\item Give an example of a complete nonflat Riemannian manifold with nonnegative sectional curvature that contains a line.
\end{enumerate}
\end{exercise}

\begin{exercise}
\label{prob:constant-sectional-curvature-immersion}
\dcref{ch12:12-1}
Suppose $(M, g)$ is a simply connected (but not necessarily complete) Riemannian $n$-manifold with constant sectional curvature. Prove that there exists an isometric immersion from $M$ into one of the model spaces $\mathbb{R}^n$, $S^n(R)$, $\mathbb{H}^n(R)$.
\end{exercise}

\begin{exercise}
\label{prob:fixed-points-orthogonal-map}
\dcref{ch12:12-2}
\uses{ex:real-projective-space-metric}
Prove that if $n$ is even, then every orientation-preserving orthogonal linear map from $\mathbb{R}^{n+1}$ to itself fixes at least one point of the unit sphere, and use this to prove that every even-dimensional spherical space form is isometric to either a round sphere or a real projective space with a constant multiple of the metric defined in Example 2.34. (Used on p. 349.)
\end{exercise}

\begin{exercise}
\label{prob:geodesic-uniqueness-nonpositive-curvature}
\dcref{ch12:12-6}
\uses{prop:path-homotopy-contains-minimizing-geodesic}
Suppose $(M,g)$ is a complete, connected Riemannian manifold, and $p,q \in M$. Proposition 6.25 showed that every path-homotopy class of paths from $p$ to $q$ contains a geodesic segment $\gamma$. Show that if $M$ has nonpositive sectional curvature, then $\gamma$ is the unique geodesic segment in the given path homotopy class.
\end{exercise}

\begin{exercise}
\label{prob:lattices-flat-tori}
\dcref{ch12:12-4}
A \emph{lattice} in $\mathbb{R}^n$ is an additive subgroup $\Lambda \subseteq \mathbb{R}^n$ of the form
\[
\Lambda = \{m_1 v_1 + \cdots + m_n v_n : m_1, \ldots, m_n \in \mathbb{Z}\},
\]
for some basis $(v_1, \ldots, v_n)$ of $\mathbb{R}^n$. Let $T^n$ denote the $n$-torus.
\begin{enumerate}
\item[(a)] Show that if $g$ is a flat Riemannian metric on $T^n$, then $(T^n, g)$ is isometric to a Riemannian quotient of the form $\mathbb{R}^n/\Lambda$ for some lattice $\Lambda$, acting on $\mathbb{R}^n$ by vector addition.
\item[(b)] Show that two such quotients $\mathbb{R}^n/\Lambda_1$ and $\mathbb{R}^n/\Lambda_2$ are isometric if and only if $\Lambda_2 = A(\Lambda_1)$ for some $A \in O(n)$.
\end{enumerate}
\end{exercise}

\begin{exercise}\label{prob:milnor-fundamental-group}\dcref{ch12:12-11}
\uses{thm:milnor-polynomial-growth,prob:nonpositive-curvature-focal-points}
Prove Theorem 12.29 (Milnor's theorem on the fundamental group of a manifold with nonnegative Ricci curvature) as follows: Let $\pi : \tilde{M} \to M$ be the universal covering space of $M$, and give $\tilde{M}$ the pullback metric $\tilde{g} = \pi^* g$.
Because the covering group is isomorphic to $\pi_1(M)$ (Prop.\ C.22), it suffices to prove the result for every finitely generated subgroup of the automorphism group. Let $\Gamma$ be such a subgroup, and let $S = \{\varphi_1, \ldots, \varphi_N\}$ be a finite generating set. Choose a base point $p \in \tilde{M}$. For each $i = 1, \ldots, N$, let $d_i$ be the Riemannian distance from $p$ to $\varphi_i(p)$, and let $D = \max(d_1, \ldots, d_n)$.
For each nonidentity element $\psi \in \Gamma$, define the \emph{length of $\psi$} to be the smallest integer $m$ such $\psi$ can be expressed as a product of $m$ elements of $S$ and their inverses.

\begin{enumerate}[label=(\alph*)]
\item Show that there exists $\varepsilon > 0$ such that for any two distinct elements $\psi_1, \psi_2 \in \Gamma$, the metric balls $B_\varepsilon(\psi_1(p))$ and $B_\varepsilon(\psi_2(p))$ are disjoint.
\item Show that if $\psi \in \Gamma$ has length $m$, then $B_\varepsilon(\psi(p)) \subseteq B_{mD+\varepsilon}(p)$.
\item Use the Bishop-Gromov theorem to prove that the number of distinct elements of $\Gamma$ of length at most $m$ is bounded by a constant times $n^m$.
\end{enumerate}
\end{exercise}

\begin{exercise}\label{prob:myers-scalar-counterexample}\dcref{ch12:12-10}
Give a counterexample to show that Myers's theorem is false in dimensions greater than 2 if we replace the lower bound on Ricci curvature by a positive lower bound on scalar curvature.
\end{exercise}

\begin{exercise}
\label{prob:no-negative-curvature-compact-products}
\dcref{ch12:12-8}
\uses{cor:no-product-negative-sectional-curvature}
Prove Corollary 12.23 (nonexistence of negatively curved metrics on compact product manifolds).
\end{exercise}

\begin{exercise}\label{prob:no-positive-negative-ricci}\dcref{ch12:12-12}
Prove that there is no compact manifold that admits both a metric of positive definite Ricci curvature and a metric of nonpositive sectional curvature.
\end{exercise}

\begin{exercise}
\label{prob:triangle-inequalities-cartan-hadamard}
\dcref{ch12:12-7}
\uses{prop:geodesic-triangle-inequalities}
Prove Proposition 12.10 (inequalities for triangles in Cartan--Hadamard manifolds). [Hint: Compare with appropriate triangles in Euclidean space.]
\end{exercise}
